\documentclass[10pt,a4paper]{amsart}
\usepackage[margin=19mm]{geometry}
\usepackage{amsmath,amssymb,mathtools}
\usepackage[scr=rsfs,cal=euler]{mathalfa}
\usepackage{xfrac}
\usepackage[unicode,hidelinks,bookmarks=false]{hyperref}
\newcommand{\ie}{i.e.\ }
\newcommand{\cf}{cf.\ }
\newcommand{\resp}{respectively\ }
\newcommand{\Subsection}{\subsection}
\providecommand{\nobreakdash}{\nobreak\hbox{-}}
\setlength{\emergencystretch}{2em}
\multlinegap0pt
\usepackage{bm}
\usepackage{bbm}
\usepackage{dsfont}
\usepackage{xspace}
\usepackage{cancel}
\usepackage{mathtools}
\usepackage{amsthm}
\makeatletter
\@ifundefined{thm}{\newtheorem{thm}{Thm}}{}
\@ifundefined{prop}{\newtheorem{prop}[thm]{Proposition}}{}
\makeatother
\makeatletter
\newcommand{\ds}{\displaystyle}
\@ifundefined{remark}{\newenvironment{remark}{\begin{preremark}\rm}{\medskip \end{preremark}}}{}
\newcommand{\vphi}{\varphi}
\makeatother
\title[Harnack inequality for non-uniformly elliptic equations in n]{Harnack inequality for non-uniformly elliptic equations in non-divergence form}
\author{David Bowman}
\date{Source: arXiv:2604.13303v1 (CC BY 4.0)}
\begin{document}
\maketitle

\noindent\emph{Source and licence.} Harnack inequality for non-uniformly elliptic equations in non-divergence form. Version arXiv:2604.13303v1 (CC BY 4.0), distributed under the Creative Commons Attribution 4.0 International licence, as recorded in the arXiv licence metadata for this exact version and shown on the abstract page https://arxiv.org/abs/2604.13303v1. Full rights notice: licenses/arxiv-2604.13303-CC-BY-4.0.txt.\par\smallskip

\noindent\textit{Editorial scope.} Counted unit: the Prop whose source label is \texttt{prop: noHarnack} in the article (source0.tex lines 392--406, source label \texttt{prop: noHarnack}), delivered together with the article's own complete original proof of it (source0.tex lines 407--493, one \texttt{proof} environment of 87 lines). No mathematics is written, repaired or re-derived: statement and proof are the article's own text line for line. Required context retained uncounted: none. Every definition, display and label of those lines is retained unchanged. Omitted from this package and not redistributed: 1--391, 494--1118. Nothing inside a retained excerpt is omitted or altered: every retained body is delivered exactly as the article's lines are written, so no omission receipt of any kind accompanies this package. Citations: the retained excerpts carry no citation command at all, so no bibliography entry of the article is transcribed and no reference is redistributed. Reference adaptation: none was needed and none was made; every reference inside the delivered unit resolves inside it, so no unresolved reference or citation is produced. Printed slips kept literal: no printed word, symbol, hypothesis or number of the counted statement or of its proof was altered. Layout and class: the article's own class is \texttt{amsart} with packages including inputenc, geometry, latexsym, amsmath, color, amsthm, epsfig, mathrsfs, mathdots, mathtools, graphicx, amssymb, mdframed, fancyhdr; the wrapper is the lane's pinned \texttt{amsart} editorial wrapper at 10pt on A4, which restates the theorem-like environments the retained text needs and adds the article's own macro definitions that the retained text needs (\texttt{\textbackslash ds}, \texttt{\textbackslash vphi}), with extra wrapper packages (bm, bbm, dsfont, xspace, cancel, mathtools). The delivered numbering is the unit's own. Assets: no retained span contains a figure environment, a graphics-inclusion command, a PDF-page inclusion, a table or a TikZ picture; no image file is copied into this package, the package declares no asset dependency and no image file is redistributed with it. Licence: version arXiv:2604.13303v1 of this article is distributed under the Creative Commons Attribution 4.0 International licence, as recorded in the arXiv licence metadata for this exact version and shown on the abstract page https://arxiv.org/abs/2604.13303v1. A full rights notice is delivered at licenses/arxiv-2604.13303-CC-BY-4.0.txt. Characters: every retained excerpt is pure ASCII, so the delivered engine, XeLaTeX with its default Latin Modern text font, reports no missing character.
\begin{prop} 
\label{prop: noHarnack} Consider any dimension $d \ge 3$ and fix an exponent $p<d-1$. Then for every $r>0$ small, there exist coefficients $a_{ij}^r$ and corresponding nonnegative viscosity solutions $u_r$ on the ball $B_1$ such that 
\begin{itemize}
\item $a_{ij}^r$ is uniformly elliptic \text{ in } $B_1$ 
\item $I \ge a_{ij}^r(x) \ge \lambda_r(x)  I$ where $\lambda_r$ is continuous and satisfies 
$$\ds \limsup_{r\to 0} \|\lambda_r^{-1}\|_{L^p(B_{1})}  <\infty,$$
\item $u_r$ is a $C^{1,1}$ viscosity solution to 
$$a_{ij}^r \partial_{ij}u_r = 0 \text{ on } B_{1}$$
\end{itemize}
and yet
$$\frac{\sup_{B_{1/2}}u_r}{\inf_{B_{1/2}}u_r} \to +\infty \text{ as } r\to 0,$$
showing that a Harnack inequality does not hold in the case $p<d-1$. In fact, the solutions $u_r$ can be chosen such that $\inf_{B_1} u_r \to 1$ as $r\to 0$, but 
$$| \{u_r \le N\}\cap B_1| \to 0 \text{ as } r \to 0$$
for any finite number $N$. Therefore, there is no possible \textit{Weak Harnack} inequality in the case $p<d-1$, either.
\end{prop}

\begin{proof}
We construct our example by smoothing out cusps in $d-1$ variables, which curve ever so slightly in the transverse direction. Fix $r>0$ small. With $p<d-1$ fixed, we first fix a parameter $\alpha \in (0,1)$ such that 
\begin{equation}\label{eq: alphaeq} (2-\alpha)p < d-1,
\end{equation}
and, with this choice of $\alpha$ being made, then fix a parameter $\sigma>0$ such that
\begin{equation}
    \label{eq: sigmaeq}
    0<\sigma < \frac{d-1 - (2-\alpha)p}{p}.
\end{equation}
We use the familiar notation $x\coloneqq (x', x_d)$, and set $\rho \coloneqq |x'|$. With this in hand, we define
$$u_{r}(x) \coloneqq 1-\delta x_d^2 + r^{-\sigma} \vphi_r(\rho),$$
where $\delta = \delta(r)>0$ is to be decided, and the function $\vphi_r(\cdot)$ is given by 
$$\vphi_r(\rho) \coloneqq \begin{cases}
    \frac\alpha2 r^{\alpha -2} \rho^2 & \text{ for }  0 \le \rho \le r,\\
    \rho^\alpha - \left(1-\frac\alpha2\right)r^\alpha & \text{ for } \rho \ge r.
\end{cases}$$
It is not hard to check that $u_r$ is $C^{1,1}$. Before we define our coefficients for which $u_r$ is a solution, we define a constant $\beta$ by fixing
\begin{equation}
\label{eq: betaeq}
    \alpha r^{-\sigma}((d-2) -\beta(1-\alpha)) = 2\delta
\end{equation}
We remark that, as $\delta = \delta(r) \to 0$, this amounts to choosing $\beta \to \frac{d-2}{1-\alpha}$, which is just a fixed number for our purposes. It is important to note that $\beta$ does not become arbitrarily small or large as $r\to 0$. 

We denote by $e_\rho$ the radial unit vector in the $x'$ variables. Now, we define the coefficients piecewise as follows:
$$A^r(x)\coloneqq \begin{cases} A_{\text{in}}(x) & \text{ for } |x'| \le r\\
A_{\text{out}}(x) & \text{ for } |x'|>r,
\end{cases}$$
where, with obvious abuses of notation,
$$A_{\text{out}}(x) \coloneqq \rho^{2-\alpha}\left(\beta e_\rho\otimes e_\rho + (I_{d-1} - e_\rho \otimes e_\rho)\right) + e_d \otimes e_d$$
and
$$A_{\text{in}}(x) \coloneqq 
\text{diag}(\lambda_{\text{in}}, \ldots, \lambda_{\text{in}}, 1),$$
with the constant $\lambda_{\text{in}}$ being given by 
$$\lambda_{\text{in}} \coloneq \ds r^{2-\alpha +\sigma}(\log 1/r)^{-1}.$$

Then, we certainly have
$$\lambda_r(x) I \le A_r(x)\le C I \text{ in } B_1,$$
where the constant $C$ depends only on $p,d,$ and $\alpha$, and
$$\lambda_r(x) \simeq \begin{cases} \lambda_{\text{in}}& \text{ for } |x'| \le r\\
(1-\frac{|x'|-r}{r})\lambda_{\text{in}} + \frac{|x'|-r}{r} |x'|^{2-\alpha} & \text{ for } r\le |x'| \le 2r,\\
|x'|^{2-\alpha} & \text{ for } |x'|>2r.
\end{cases}$$
The only point defining $\lambda_r(\cdot)$ on the annulus $r\le |x'| \le 2r$ in this way was to demonstrate that the minimal eigenvalue of the matrix $A$ can be minorized by a \textit{continuous} function satisfying the requisite $L^p$ bound. 

Integrating in cylindrical coordinates, one finds 
\begin{align} \int_{B_1} \lambda_r^{-p} &\lesssim_d \int_{|x'| \le 2r} r^{p(\alpha-2) - p\sigma}(\log(1/r))^p \ dx' + \int_{2r <|x'| <1} |x'|^{p(\alpha-2)} \ dx'\\
& \simeq_{d,p,\alpha} r^{p(\alpha-2) - p\sigma+ d-1}(\log 1/r)^{p} + 1
\end{align}
which stays bounded as $r\to 0$ precisely under the stated conditions \eqref{eq: alphaeq} and \eqref{eq: sigmaeq} on $\alpha$ and $\sigma$. 

Next, we check that $u_r$ is a viscosity solution on $B_1$, and in fact classical away from the interface $|x'|=r$. Indeed, for $|x'| <  r$, we simply have
$$a_{ij}^r \partial_{ij}u_r = \alpha  (d-1)\lambda_{\text{in}} r^{\alpha-2-\sigma} - 2 \delta=0$$
provided we choose
$$\delta = \delta(r) \coloneqq \frac{\alpha(d-1)}{2} (\log 1/r)^{-1}, $$
which certainly tends to $0$ as $r\to 0$. Next, we remark that that the map $x' \mapsto |x'|^\alpha$ has radial curvature equal to  $\alpha(\alpha -1) \rho^{\alpha-2}$ and $d-2$ tangential curvatures equal to $\alpha \rho^{\alpha-2}$. Then we may see that on $|x'|>r$, there holds
\begin{align} a_{ij}^r \partial_{ij}u_r &= \rho^{2-\alpha}r^{-\sigma}\left(\beta \alpha(\alpha-1) \rho^{\alpha -2} + \alpha(d-2) \rho^{\alpha -2}\right) - 2\delta\\
&= \alpha r^{-\sigma}(d-2 - \beta(1-\alpha)) - 2\delta =0
\end{align}
precisely under the stated condition \eqref{eq: betaeq}.

Finally, we briefly explain how to see that $u$ satisfies the definition of viscosity solution on the sphere $\{|x'|=r\}$. Fix such an $x_0$. We work in the cylindrical frame $e_\rho, \tau_1, \ldots \tau_{d-2}, e_d$, where $\tau_1, \ldots \tau_{d-2}$ are tangent to the sphere $\{|x'|=r\}$. One may check that $\partial_{\tau_i \tau_i} u_r$ and $\partial_{dd}u_r$ are continuous across the sphere, taking values
$$\partial_{\tau_i \tau_i} u_r = a \coloneqq \alpha r^{\alpha -2-\sigma}$$
$$\partial_{dd}u_r = -2\delta.$$
On the other hand, $\partial_{\rho \rho} u_r$ has a jump across the sphere. Indeed, 
$$\lim_{s \to 0^+} \partial_{\rho\rho}u(x_0 - se_\rho) = a$$
while
$$\lim_{s \to 0^+} \partial_{\rho\rho}u(x_0+ se_\rho) = b \coloneqq \alpha(\alpha-1)r^{\alpha-2-\sigma} <0<a.$$

Accordingly, if $\psi$ touches $u_r$ from above at $x_0 \in \{|x'|=r\}$, then, restricting $u_r - \psi$ to the inward normal half-line, the tangential lines, and the $x_d$-line, one sees that
$$\partial_{\rho\rho}\psi(x_0) \ge a, \qquad \partial_{\tau_i\tau_i} \psi(x_0) \ge a, \qquad \partial_{dd} \psi(x_0) \ge - 2\delta.$$
Thus
$$a_{ij}^r(x_0)\partial_{ij}\psi(x_0) = \lambda_{\text{in}}\left(\partial_{\rho \rho}\psi(x_0) + \sum_{i=1}^{d-2} \partial_{\tau_i\tau_i}\psi(x_0)\right) + \partial_{dd}\psi(x_0) \ge (d-1)\lambda_{\text{in}}a - 2\delta =0.$$

On the other hand, if $\psi$ touches $u_r$ from below at $x_0 \in \{|x'|=r\}$, then we restrict $u_r -\psi$ to the outward normal half-line and use otherwise analogous reasoning to find 
$$\partial_{\rho\rho}\psi(x_0) \le b, \qquad \partial_{\tau_i\tau_i} \psi(x_0) \le a, \qquad \partial_{dd} \psi(x_0) \le - 2\delta.$$
Then
$$a_{ij}^r(x_0) \partial_{ij}\psi(x_0) \le \lambda_{\text{in}}(b+ (d-2)a) -2\delta \le 0$$
since $b<a$. We conclude that $u_r$ is indeed a viscosity solution on the whole ball.


We conclude the proof by noticing that
$$\inf_{B_{1/2}} u_r = 1- \frac{\delta(r)}{4} \to 1 \text{ as } r\to 0,$$
while
$$\sup_{B_{1/2}}u_r = 1 + r^{-\sigma}\left(2^{-\alpha} -\left(1-\frac\alpha2\right)r^\alpha\right) \to +\infty \text{ as } r\to 0.$$
Furthermore, it is clear that $|\{u_r \le N\} \cap B_1|\to 0$ as $r\to 0$ for any finite $N<\infty$; the functions $u_r$ converge locally uniformly to $+\infty$ away from the axis $\{x'=0\}$. Indeed, one may compute that
$\{u_r \le N\} \cap B_1 \subset \{(x',x_d)\subset B_1 \ : \ |x'| \le C_N r^{\sigma/\alpha}\}$, assuming we also choose $\sigma<\alpha$. 
\end{proof} 

\end{document}
